% ЕМТ 2023, X клас — точен LaTeX препис от официалните файлове в Klasirane.com.
% Условия: https://klasirane.com/api/competitions/EMT/All/2023/10%20%D0%BA%D0%BB/probs
% SHA-256 (условия): e6ae93a8265af3c0b3b7dbcfe542c3c0685ec5ea6126ce1df9e716f48d9a5f3b
% Решения: https://klasirane.com/api/competitions/EMT/All/2023/10%20%D0%BA%D0%BB/sol
% SHA-256 (решения): 72280a310fecfa9629091c07e18403b28b1eaf5ef75e16459ae36c6e8706e540
% Точковите оценки, правописът, формулировките и видимите печатни особености са запазени от източника.
% В официалните файлове няма отделна чертежна фигура; не е добавяна такава.

Задача 10.1. Да се реши уравнението
$$
(x+1)\sqrt{x^2+2x+2}+x\sqrt{x^2+1}=0.
$$

Решение. Отговор: $x=-\dfrac12$.

Първи начин: Записваме уравнението във вида $(x+1)\sqrt{x^2+2x+2}=-x\sqrt{x^2+1}$, забелязваме, че и двата израза под корен са строго положителни за всяко реално $x$, като освен това
$$
-x(x+1)\geq0,
$$
откъдето $x\in[-1,0]$. Повдигаме на втора степен двете страни и преработваме:
$$
\begin{aligned}
(x^2+2x+1)(x^2+2x+2)&=x^2(x^2+1)\\
x^4+4x^3+7x^2+6x+2&=x^4+x^2\\
2x^3+3x^2+3x+1&=0\\
(2x+1)(x^2+x+1)&=0.
\end{aligned}
$$
Тъй като $x^2+x+1>0$, то единственото възможно решение е $x=-\dfrac12$. Тъй като $-\dfrac12\in[-1,0]$, то $x=-\dfrac12$ действително е решение.

Втори начин: Полагаме $a=\sqrt{x^2+2x+2}>0$ и $b=\sqrt{x^2+1}>0$. Следователно
$$
x=\frac{(x^2+2x+2)-(x^2+1)-1}{2}=\frac{a^2-b^2-1}{2}
$$
и
$$
x+1=\frac{a^2-b^2+1}{2}.
$$
Уравнението е еквивалентно на:
$$
\begin{aligned}
\frac{a^2-b^2+1}{2}\cdot a+\frac{a^2-b^2-1}{2}\cdot b&=0\\
(a^2-b^2)a+a+(a^2-b^2)b-b&=0\\
(a-b)((a+b)^2+1)&=0.
\end{aligned}
$$
Оттук, $a=b$ и $x=-\dfrac12$.

Оценяване. (6 точки) 2т. за $x\in[-1,0]$; 2т. за извеждане на уравнението $2x^3+3x^2+3x+1=0$; 2т. за $x=-\dfrac12\in[-1,0]$. Ако е повдигнато на втора степен без да са направени ограничения, задачата се оценява най-много на 5т. (ако има проверка за $x=-\dfrac12$).

Алтернативно: (6 точки) 2т. за полагане на $a,b$; 2т. за изразяване на $x$ и $x+1$; 2т. за $a=b$ и $x=-\dfrac12$.

Задача 10.2. Даден е остроъгълен триъгълник $ABC$ с център $O$ на описаната окръжност. Точка $P$ е върху страната $BC$, такава че $BP<\dfrac12BC$. Точка $Q$ е от страната $BC$ такава, че $CQ=BP$. Правата $AO$ пресича $BC$ в точка $D$, а точка $N$ е среда на $AP$. Описаната около триъгълник $ODQ$ окръжност пресича за втори път описаната около триъгълник $BCO$ окръжност в точка $E$. Правите $NO$ и $OE$ пресичат $BC$ съответно в точки $K$ и $F$. Да се докаже, че точките $A,O,K,F$ лежат на една окръжност.

Решение. Нека $AO$ пресича описаната около $\triangle ABC$ окръжност в точка $L$, а описаната около $\triangle BOC$ окръжност в точка $R$. Да означим с $X$ средата на $BC$, а с $Y$ – диаметрално противоположната на $O$ в описаната около $\triangle BOC$ окръжност. Нека $OQ$ пресича за втори път описаната около $\triangle BOC$ окръжност в точка $T$.

Имаме $OX\cdot OY=OD\cdot OR=OQ\cdot OT=OE\cdot OF$. От последните три равенства лесно следва, че $R,T,F$ лежат на една права. От $\angle ORF=\angle OQP=\angle QPO=\angle FPO$ следва, че $OPRF$ е вписан и значи $PD\cdot DF=OD\cdot DR=DL\cdot DA$. Оттук, $APLF$ е вписан. Но $PL\parallel NO$ (тъй като $NO$ е средна отсечка в $\triangle PLA$) и значи $AOKF$ също е вписан.

Оценяване. (6 точки) 2т. за $OX\cdot OY=OD\cdot OR=OQ\cdot OT=OE\cdot OF$; 3т. за $APLF$ вписан; 1т. за довършване.

Задача 10.3. Да се намерят всички естествени числа $k$ със следното свойство: Съществува полином $f(x)$ с рационални коефициенти, такъв че за всяко естествено число $n>2023^{2023}$
$$
f(n)=\operatorname{НОК}(n+1,n+2,\ldots,n+k).
$$

Решение. При $k=1$ и $k=2$ търсените полиноми са съответно $f(x)=x+1$ и $f(x)=(x+1)(x+2)$. Нека $k\geq3$ и да допуснем, че съществува такъв полином $f(x)$. За всяко просто число $p$, степента му в $\operatorname{НОК}(n+1,n+2,\ldots,n+k)$ е $\max\{\alpha_1,\alpha_2,\ldots,\alpha_k\}$, където $\alpha_i$ е степента на $p$ в каноничното представяне на $n+i$, $i=1,\ldots,k$. Ако това е примерно $\alpha_s$, то тя би се получила ако вземем
$$
\frac{(n+1)(n+2)\cdots(n+k)}{p^{\alpha_1}p^{\alpha_2}\cdots p^{\alpha_{s-1}}p^{\alpha_{s+1}}\cdots p^{\alpha_k}},
$$
като е ясно, че степените на $p$ в знаменателя са делители на $\displaystyle\prod_{1\leq i\ne s\leq k}(s-i)$. Следователно
$$
\operatorname{НОК}(n+1,n+2,\ldots,n+k)=\frac{(n+1)(n+2)\cdots(n+k)}{C_n}, \tag{1}
$$
където $C_n$ е делител на $\displaystyle\prod_{1\leq i<j\leq k}(j-i)$. Понеже $C_n$ може да приема краен брой допустими стойности, ще има естествено число $C$ такова, че за безброй много $n$, $f(n)=\dfrac{(n+1)(n+2)\cdots(n+k)}{C}$. Значи за безброй много $x$ $f(x)=\dfrac{(x+1)(x+2)\cdots(x+k)}{C}$, откъдето
$$
f(x)=\frac{(x+1)(x+2)\cdots(x+k)}{C},\qquad \forall x\in\mathbb R.
$$
Следователно
$$
\operatorname{НОК}(n+1,n+2,\ldots,n+k)=\frac{(n+1)(n+2)\cdots(n+k)}{C},\qquad \forall n\in\mathbb N.
$$
Да допуснем, че това е възможно. Да изберем просто число $p<k$ такова, че $p$ не дели $k$. Нека $n+k+1=p^m$ за достатъчно голямо $m$. От горната формула имаме
$$
\frac{\operatorname{НОК}(n+2,n+3,\ldots,n+k+1)}{\operatorname{НОК}(n+1,n+2,\ldots,n+k)}=\frac{n+k+1}{n+1}. \tag{2}
$$
Степента на $p$ в числителя на лявата страна е $m$, а в знаменателя – поне 1, докато степента на $p$ в числителя на дясната страна е $m$, а в знаменателя – 0. Противоречие! Следователно, допускането е грешно и при $k\geq3$ не съществува полином с исканото свойство.

Оценяване. (7 точки) 1т. за случая $k=2$; 3т. за (1); 1т. за съществуването на $C$; 2т. за извеждане на (2) и избор на подходящо $n$.

Задача 10.4. Във всяка клетка на таблица $101\times101$ е записано естествено число. Известно е, че както и да изберем 101 клетки на таблицата, никои две от които не лежат в един ред или стълб, сумата на числата в избраните клетки се дели на 101. Да се докаже, че броят начини да изберем по една клетка от всеки ред на таблицата така, че сумата на числата в избраните клетки да се дели на 101, се дели на 101.

Решение. Нека $p=101$ и да номерираме редовете и стълбовете на таблицата с числата от 1 до $p$. С $c_{r,s}$ ще означаваме числото, записано в ред $r$ и стълб $s$ на таблицата. Ще наричаме ключалка множество от $p$ клетки, никои две от които не лежат в един ред или стълб. Нека $i\ne j,k\ne l\in\{1,2,\ldots,p\}$ - лесно се вижда, че можем да изберем $p-2$ клетки, които допълват двойките клетки $(i,k),(j,l)$ и $(i,l),(j,k)$ до ключалки. Тогава от условието следва, че
$$
c_{i,k}+c_{j,l}\equiv c_{i,l}+c_{j,k}\pmod p.
$$
Нека сега изберем цели числа $a_1,a_2,\ldots,a_p,b_1,b_2,\ldots,b_p$ такива, че за всяко $i$ да е изпълнено $c_{1,i}=a_1+b_i$ и $c_{i,1}=a_i+b_1$. Тогава лесно се проверява, че за всеки $r,s$ е изпълнено
$$
c_{r,s}\equiv a_r+b_s\pmod p.
$$
Ще наричаме обобщена ключалка множество от $p$ клетки, никои две от които не се намират в един и същ ред. Искаме да докажем, че броят на обобщените ключалки със сума, кратна на $p$, е кратен на $p$. Да разгледаме клетките $(1,k_1),(2,k_2),\ldots,(p,k_p)$ и нека с $d_j$ означим броя на тези клетки, които се намират в стълб $j$. Имаме:
$$
\sum_{i=1}^p c_{i,k_i}\equiv\sum_{i=1}^p(a_i+b_{k_i})\equiv\sum_{i=1}^p a_i+\sum_{i=1}^p d_i b_i\pmod p.
$$
Това показва, че остатъкът $\pmod p$ на сумата на числата в дадена обобщена ключалка зависи единствено от набора $(d_1,d_2,\ldots,d_p)$. Да разгледаме набор, за който по-горната сума е кратна на $p$. Броят обобщени ключалки, които имат този набор, е
$$
\frac{p!}{\prod_{i=1}^p d_i!}. \tag{1}
$$
Това число се дели на $p$, стига поне две от числата $d_1,d_2,\ldots,d_p$ да са ненулеви (тъй като $p$ е просто). Случаят, в който точно едно от тези числа е ненулево, съответства на обобщена ключалка, в която всички клетки са в един и същ стълб. Следователно броят обобщени ключалки, различни от стълб на таблицата, със сума, кратна на $p$, се дели на $p$. Остава да забележим, че за сумата на числата в $i$-тия стълб на таблицата имаме
$$
\sum_{j=1}^p c_{j,i}\equiv\sum_{j=1}^p(a_j+b_i)\equiv\sum_{j=1}^p a_j\pmod p,
$$
и тогава или всички стълбове имат сума, кратна на $p$, или нито един от стълбовете не е с такава сума. Исканото следва.

Забележка. Твърдението на задачата остава вярно за произволна таблица $101\times101$ (и съответно $p\times p$, където $p$ е просто число).

Оценяване. (7 точки) 1т. за $c_{i,k}+c_{j,l}\equiv c_{i,l}+c_{j,k}\pmod p$; 1т. за $c_{r,s}\equiv a_r+b_s\pmod p$; 1т. за наблюдението, че сумата на числата в дадена обобщена ключалка зависи единствено от набора $(d_1,d_2,\ldots,d_p)$; 2т. за (1); 2т. за довършване.
